\section{CSR, Privilege, Exceptions, and Traps}

\subsection{Why CSRs exist}

General-purpose registers hold normal program data. A CPU also needs control/status registers for privileged state: interrupt enables, trap handler address, trap cause, saved exception PC, timers, and address translation control.

RISC-V calls these \textbf{CSRs}: Control and Status Registers.

\subsection{What a trap is}

A trap is a controlled redirection of execution caused by an exception or interrupt.

\begin{itemize}
\item An \textbf{exception} is caused by the current instruction, such as illegal instruction or \texttt{ecall}.
\item An \textbf{interrupt} is caused by an external/asynchronous event, such as a timer.
\end{itemize}
On a trap, the CPU must save enough information so software can handle the event and later return.

Conceptually:

\begin{lstlisting}[language={}]
mepc   = PC of interrupted/faulting instruction
mcause = reason for trap
PC     = mtvec
mode   = machine mode
\end{lstlisting}

\subsection{Privilege model in this design}

\texttt{rtl/csr.v} implements a minimal privilege model:

\begin{longtable}{|>{\RaggedRight\arraybackslash}p{0.28\textwidth}|>{\RaggedRight\arraybackslash}p{0.32\textwidth}|>{\RaggedRight\arraybackslash}p{0.30\textwidth}|}
\hline
\textbf{Mode} & \textbf{Encoding} & \textbf{Meaning} \\
\hline
\endfirsthead
\hline
\textbf{Mode} & \textbf{Encoding} & \textbf{Meaning} \\
\hline
\endhead
U-mode & \texttt{2'b00} & User mode, restricted. \\
\hline
M-mode & \texttt{2'b11} & Machine mode, privileged. \\
\hline
\end{longtable}

The design resets into machine mode. Traps enter machine mode. \texttt{mret} restores the previous privilege mode using bits saved in \texttt{mstatus}.

\subsection{Supported CSRs}

\texttt{rtl/csr.v} supports:

\begin{longtable}{|>{\RaggedRight\arraybackslash}p{0.28\textwidth}|>{\RaggedRight\arraybackslash}p{0.32\textwidth}|>{\RaggedRight\arraybackslash}p{0.30\textwidth}|}
\hline
\textbf{CSR} & \textbf{Address} & \textbf{Role} \\
\hline
\endfirsthead
\hline
\textbf{CSR} & \textbf{Address} & \textbf{Role} \\
\hline
\endhead
\texttt{mstatus} & \texttt{0x300} & Global interrupt enable and previous interrupt enable bits. \\
\hline
\texttt{mie} & \texttt{0x304} & Machine interrupt enable bits, including timer/external. \\
\hline
\texttt{mtvec} & \texttt{0x305} & Trap handler base address. \\
\hline
\texttt{mscratch} & \texttt{0x340} & Scratch register for trap handlers. \\
\hline
\texttt{mepc} & \texttt{0x341} & Saved PC for trap return. \\
\hline
\texttt{mcause} & \texttt{0x342} & Trap cause. \\
\hline
\texttt{mip} & \texttt{0x344} & Pending interrupt bits, read-only style here. \\
\hline
\texttt{satp} & \texttt{0x180} & Address translation control for MMU variants. \\
\hline
\end{longtable}

\subsection{CSR instructions}

The CSR unit supports the CSR read/write family through the \texttt{csr\_next} function:

\begin{longtable}{|>{\RaggedRight\arraybackslash}p{0.38\textwidth}|>{\RaggedRight\arraybackslash}p{0.52\textwidth}|}
\hline
\textbf{Operation} & \textbf{Meaning} \\
\hline
\endfirsthead
\hline
\textbf{Operation} & \textbf{Meaning} \\
\hline
\endhead
\texttt{csrrw} / \texttt{csrrwi} & Write source value into CSR. \\
\hline
\texttt{csrrs} / \texttt{csrrsi} & Set CSR bits selected by source. \\
\hline
\texttt{csrrc} / \texttt{csrrci} & Clear CSR bits selected by source. \\
\hline
\end{longtable}

The CPU writes the old CSR value to \texttt{rd} and updates the CSR according to the instruction.

\subsection{10.6 Trap entry and \texttt{mret}}

\texttt{csr.v} prioritizes updates as:

\begin{lstlisting}[language={}]
reset > trap entry > mret > CSR write
\end{lstlisting}

On trap entry:

\begin{itemize}
\item \texttt{mepc} receives the trap PC;
\item \texttt{mcause} receives the cause;
\item \texttt{mstatus.MPIE} receives old \texttt{MIE};
\item \texttt{mstatus.MIE} is cleared to disable nested interrupts by default;
\item current privilege is saved into \texttt{mstatus.MPP};
\item current privilege becomes machine mode.
\end{itemize}
On \texttt{mret}:

\begin{itemize}
\item \texttt{MIE} is restored from \texttt{MPIE};
\item privilege is restored from \texttt{MPP};
\item \texttt{MPP} is cleared down toward user mode;
\item PC redirects to \texttt{mepc}.
\end{itemize}
\subsection{Illegal instruction detection}

Trap-capable CPU wrappers compute whether the opcode/instruction is known. If not, they raise illegal instruction cause \texttt{2}. They also treat user-mode attempts to execute privileged instructions such as CSR access or \texttt{mret} as privilege violations.

\bigskip
